\section{关系与等价类}
\label{sec:01-Mathematical-Language/02-Sets-Relations-Functions/02}

\subsection{从一个日常问题出发}

假设你有三个学生：阿明、小华、大伟，还有三门课：数学、物理、编程。你手上有一张选课表：阿明选了数学和编程，小华选了物理，大伟选了数学和物理。

怎么把这张表写成数学语言？

一种自然的想法：把所有可能的``谁选了什么课''列出来。三个学生、三门课，可能的组合一共 \(3\times 3 = 9\) 种：

\[
\{(\text{阿明},\text{数学}),\;(\text{阿明},\text{物理}),\;(\text{阿明},\text{编程}),\;
 (\text{小华},\text{数学}),\;\dots,\;(\text{大伟},\text{编程})\}.
\]

选课表做的事情，就是从这 9 个可能性里挑出一部分打上勾——挑出来的就是实际选课记录：

\[
\{(\text{阿明},\text{数学}),\;(\text{阿明},\text{编程}),\;
  (\text{小华},\text{物理}),\;(\text{大伟},\text{数学}),\;(\text{大伟},\text{物理})\}.
\]

你看，选课关系不过是从 9 个候选有序对中选出了 5 个。

这就是``关系''的全部秘密。

\subsection{关系就是一堆有序对}

从集合 \(A\) 到集合 \(B\) 的二元关系 \(R\)，是笛卡尔积 \(A\times B\) 的一个子集。如果 \((a,b)\in R\)，写成 \(aRb\)。

笛卡尔积提供所有可能的配对，关系做的事情就是回答一个简单问题：哪几对算数？

整数上的整除关系是 \(\Z\times\Z\) 的一个子集——\((6,3)\) 在里面（因为 \(6\) 能整除 \(3\)？不对，\(3\) 能整除 \(6\)，所以 \((3,6)\) 在里面，\((6,3)\) 不在）；实数上的``小于''关系也是 \(\R\times\R\) 的子集。画成箭头、列成表格、写成条件语句，都是同一堆有序对的不同脸孔。

当关系是从 \(A\) 到自身的（也就是 \(A\times A\) 的子集），我们就可以问三个分类性的问题：

\begin{itemize}
\tightlist
\item
  自反性：每个 \(a\in A\) 都跟自己有关系吗？也就是 \((a,a)\) 在不在 \(R\) 里；
\item
  对称性：只要 \((a,b)\) 在，\((b,a)\) 就一定在吗？
\item
  传递性：\((a,b)\) 和 \((b,c)\) 都在，\((a,c)\) 就一定也跟着在吗？
\end{itemize}

这三条是彼此独立的——满足其中一条不承诺另外两条。你看：

``认识''这个关系，大体对称（你认识我我也认识你），但常常不传递——你认识小华，小华认识小华的表哥，不代表你也认识那位表哥。

``小于''呢？传递得很（\(2<3\) 且 \(3<5\) 推出 \(2<5\)），但它既不自反（\(2<2\) 是假的），也不对称（\(2<3\) 成立但 \(3<2\) 不成立）。

``等于''最乖，三条全满足。

所以每次遇到一个关系，别凭直觉猜它有什么性质。拿具体元素一条一条过。比如 \(a=3,b=5\) 能验证传递性吗？不能——需要三个点才能验证，一个例子不够。

\subsection{等价关系：故意忽略一些差异}

同时满足自反、对称、传递的关系，叫等价关系。

它不代表两个东西完全相同。它只是说：按当前关心的标准，它们可以归成一类。

想一想你是怎么办身份证的。公安机关不关心你今天的发型和十年前是不是一样，不关心你昨天吃了什么。它只锁定几个字段——姓名、出生日期、身份证号——如果这些吻合，两笔记录就判为同一个人。你把无数种差异全部忽略，只盯着选中的特征。等价关系就是这种``锁定少数特征，其余全部忽略''的数学装置。

来看一个你迟早会反复碰到的例子。固定一个正整数 \(m\)，在整数上定义：

\[
a\sim b
\quad\Longleftrightarrow\quad
m\mid(a-b).
\]

读作：\(a\) 和 \(b\) 除以 \(m\) 的余数相同。\(m\mid(a-b)\) 是``\(m\) 整除 \(a-b\)''，即 \(a-b\) 是 \(m\) 的倍数。

验证一下三条性质。自反性：\(a-a=0\) 被任何 \(m\) 整除，所以 \(a\sim a\)；对称性：\((a-b)\) 是 \(m\) 的倍数，\((b-a)\) 显然也是；传递性：\((a-b)\) 和 \((b-c)\) 都是 \(m\) 的倍数，加起来 \((a-c)\) 也是。

全通过。所以 \(\sim\) 是一个等价关系。

\subsection{等价类：关系画出的分界线}

对任意元素 \(a\)，它的等价类是：

\[
[a]=\{x\in A:x\sim a\}.
\]

就是你跟谁等价，谁就跟你待在一个房间里。

拿 \(m=3\) 实际看看。整数按除以 3 的余数分成三间房：

\[
[0]=\{\ldots,-6,-3,0,3,6,9,\ldots\},\quad
[1]=\{\ldots,-5,-2,1,4,7,\ldots\},\quad
[2]=\{\ldots,-4,-1,2,5,8,\ldots\}.
\]

注意一个极易踩坑的细节：\([1]\) 和 \([4]\) 不是两个有交叠的类——它们完全就是同一个类，因为 \(1\sim4\)（除以 3 都余 1）。等价类不关心你用哪个代表元素去叫它。类是一个整体，名字可以换。

\subsection{为什么等价类只能相等或不交}

假设 \([a]\) 和 \([b]\) 沾上了一点边——存在一个 \(x\)，既在 \([a]\) 又在 \([b]\) 里。那么 \(x\sim a\) 且 \(x\sim b\)。

由对称性：\(a\sim x\)。再配合 \(x\sim b\) 用传递性：\(a\sim b\)。

现在，任取 \([a]\) 里的一个元素 \(y\)，有 \(y\sim a\) 且 \(a\sim b\)，传递过去 \(y\sim b\)，所以 \(y\in[b]\)。这就说明了 \([a]\subseteq[b]\)。反过来倒腾一遍，\([b]\subseteq[a]\)。两个类完全重合。

推出来的结论很干脆：

\begin{enumerate}
\tightlist
\item
  每个元素都稳稳待在自己的等价类里（自反性保证）；
\item
  两个等价类，要么一模一样，要么一个共同元素都没有；
\item
  所有等价类拼起来不多不少就是原集合。
\end{enumerate}

这三条放到一起，就是集合的一个\textbf{划分}——把集合切成互不重叠的若干块。反方向也成立：你先随便给集合切块，然后宣布``在同一块里的元素等价''，自动得到一个等价关系。

``等价关系''和``把集合分组''说的是同一个结构的两面。你在不同数学分支里会反复看到这个二象性。

\subsection{商集：把类当成新元素}

所有等价类凑成的集合叫商集，记作 \(A/{\sim}\)。商集的元素不再是一个个的个体，而是整个整个的类。

有理数的构造就是一次漂亮的商集操作。从整数对 \((p,q)\)（\(q\neq 0\)）出发，规定

\[
(p,q)\sim(r,s) \quad\Longleftrightarrow\quad ps=qr.
\]

这样 \(\frac{1}{2}\) 和 \(\frac{2}{4}\) 和 \(\frac{-3}{-6}\)——三对不同的整数对——全掉进同一个等价类。商集 \(\Z\times(\Z\setminus\{0\})/{\sim}\) 就是有理数集 \(\mathbb{Q}\)。

留意为何用乘法 \(ps=qr\) 而不是除法 \(p/q=r/s\)：除法会预先假设``比值''这个东西已经存在，而我们正在构造它。等价关系让你声明``这些不同的写法代表同一个东西''，而不需要那个东西先在场。构造新数学对象时，等价关系就是干这件事的。

\subsection{传递性是埋伏最多的那条}

在实数上定义一个关系：\(xRy\) 当且仅当 \(|x-y|\le 1\)。

自反？没问题，\(|x-x|=0\le1\)。对称？绝对值天生对称。传递性呢？

试三个数：\(0,\,0.8,\,1.6\)。\(|0-0.8|=0.8\le1\)，所以 \(0R0.8\)；\(|0.8-1.6|=0.8\le1\)，所以 \(0.8R1.6\)。但 \(|0-1.6|=1.6>1\)，\(0\) 和 \(1.6\) 没关系。

传递性碎掉了。距离近不等于等价——一段段短跳板可以架到很远的地方去。很多人在自以为``近似相等''的场景里踩中这个坑。

\subsection{闭包：在原关系上打补丁}

你拿到的关系差了一条性质，又不想删掉已经确认的有序对。怎么办？硬补。缺什么补什么，补到它满足为止。

如果在所有包含 \(R\) 且具备指定性质的关系里，存在一个最小的（按集合包含排序），就叫它 \(R\) 关于该性质的\textbf{闭包}。``最小''的意思是：除了为了达标不得不加的对，一个多余的也不放进去。

举个最小的例子。集合里只有三个元素 \(a,b,c\)，原关系很穷：

\[
R=\{(a,b),\;(b,c)\}.
\]

它显然不传递——有 \(aRb\) 和 \(bRc\)，却缺了 \(aRc\)。补进去：

\[
R^+=\{(a,b),\;(b,c),\;(a,c)\}.
\]

这就是传递闭包。逻辑很清楚：只要存在一条关系链

\[
x=x_0Rx_1R\cdots Rx_n=y,
\]

传递闭包里就强行令 \(xR^+y\)。把每个元素想成一个点，每个有序对想成一条有向边，那么原关系记录直接相邻的边，传递闭包记录``能不能沿一条或多条边走通''。因此可以写作

\[
R^+=\bigcup_{n\ge 1}R^n,
\]

其中 \(R^n\) 是走恰好 \(n\) 步能连上的有序对。如果再把每个 \((x,x)\) 也塞进去，就是自反传递闭包 \(R^*\)。

同样，自反闭包补齐所有 \((x,x)\)；对称闭包看到 \((x,y)\) 就补 \((y,x)\)。闭包不是一个单一操作——你必须明确说清楚要补齐哪条性质。

\subsection{别迷信闭包}

闭包只在形式上补齐了规则。它不负责保证补出来的有序对仍然符合你心里的直觉含义。

把前面``距离不超过 1''的关系做传递闭包：\(0\) 连到 \(0.8\)，\(0.8\) 连到 \(1.6\)，\(1.6\) 连到 \(2.4\)……一路走下去，\(0\) 最终能连到任何一个实数。但这个新关系已经跟日常所说的``彼此接近''没有任何关系了。

客户记录去重是同一种陷阱。你规定``姓名和地址相似度超过 0.85''就算同一个人。记录 A 像 B，B 像 C，C 像 D……一串桥接之后，A 和 D 可能已经毫无相似之处。直接把相似关系的传递闭包当身份分组，会让一条条局部匹配把不同的人逐步并到一起。要得到真正的客户身份分组，必须引入稳定标识、可核查的合并规则，或者明确接受桥接的后果。相似关系自身不会自动进化成等价关系。

类似的，``拥有共同的朋友''也不传递。看到``相似''``接近''``相关''这些词时，别想当然地当等价关系处理——三条性质，拿具体数值，一条一条验。

延伸阅读：MIT Mathematics for Computer Science 的 Relations 章节。
